\section{Vista de contenedores}

\subsection{Diagrama de Contenedores — Kit Pagos Colombia (Nivel 2)}
\begin{figure}[htbp]
\centering
\includegraphics[width=\textwidth]{"../../architecture/SDK/Container Diagram - C4.png"}
\caption{Diagrama de Contenedores de Kit Pagos Colombia}
\label{fig:container_system}
\end{figure}

El diagrama de contenedores descompone el sistema Kit Pagos Colombia en sus unidades desplegables y ejecutables, mostrando cómo se distribuyen las responsabilidades entre ellas y cómo se comunican entre sí y con los sistemas externos.

\subsubsection{Contenedores del sistema}
El sistema está compuesto por tres contenedores principales que operan de forma complementaria.

El SDK Kit Pagos Colombia es el componente central del framework. Está desarrollado en TypeScript y se distribuye como un paquete npm compatible con Node.js 20.19 o superior. Su función principal es ofrecer una interfaz unificada para que los desarrolladores puedan crear pagos, consultar el estado de las transacciones y validar webhooks, sin tener que implementar de manera independiente la integración con cada pasarela. Internamente, el SDK utiliza una arquitectura hexagonal y cuenta con cuatro adaptadores, uno para cada pasarela soportada: Wompi, Rapyd, Mercado Pago y Kushki. Cada adaptador se encarga de transformar el contrato definido por el SDK al formato y las reglas específicas de la pasarela correspondiente. Para utilizarlo, el desarrollador instala el paquete mediante npm y configura las credenciales de las pasarelas que necesita. Cuando se trabaja en modo simulación, los adaptadores envían las solicitudes a la API de Simulación en lugar de comunicarse con las APIs reales de los proveedores. El paquete incluye además el subpaquete independiente \texttt{kit-pagos-colombia/browser}, que tokeniza tarjetas en el navegador contra Wompi, Mercado Pago y Kushki sin que el número de la tarjeta pase por el servidor del comercio.

Los cuatro adaptadores —Wompi, Rapyd, Mercado Pago y Kushki— forman parte del mismo paquete npm y están implementados como módulos internos del SDK. Por esta razón, no se consideran contenedores independientes: no tienen un proceso propio ni pueden desplegarse de manera aislada. La responsabilidad específica de cada uno se detalla posteriormente en el diagrama de componentes de Nivel 3.

La API de Simulación es un servicio REST desarrollado con Fastify y Node.js 20.19 o superior, desplegado en Render y disponible mediante una URL pública HTTPS. Su objetivo es simular el comportamiento de las cuatro pasarelas para los principales escenarios utilizados durante las pruebas, evitando la necesidad de utilizar credenciales o realizar transacciones reales. El simulador expone dos superficies HTTP bajo el mismo proceso, puerto y despliegue: \texttt{/v1/sim} imita las APIs nativas de Wompi, Rapyd, Mercado Pago y Kushki, y \texttt{/v1/api} expone las operaciones de Kit Pagos a través del SDK. La API REST no reemplaza las rutas de simulación: ambas superficies coexisten y tienen prefijos y módulos separados.

El escenario de \texttt{/v1/sim} se resuelve con una precedencia explícita: la cabecera \texttt{x-simulator-scenario} (con \texttt{x-simulate-scenario} como alias), el dato de prueba nativo de la pasarela (tarjeta, banco PSE, titular o documento), un monto reservado y, por omisión, \texttt{APPROVED}. Los escenarios cubren éxito, rechazo, pendiente, expirado, pago duplicado, límite de tasa, errores de servidor y de red, timeout, respuesta lenta y cuerpos malformados. El simulador también permite generar webhooks sintéticos firmados hacia la URL configurada por el comercio, mediante una invocación manual o de forma automática cuando cambia el estado de un pago simulado. Estos webhooks reproducen el formato y el algoritmo de firma de cada proveedor. El Tester o equipo de QA interactúa directamente con este servicio mediante HTTPS para validar los diferentes flujos de pago dentro de un entorno controlado. La documentación de la API se encuentra disponible mediante OpenAPI 3.0 en el endpoint \texttt{/docs}, con el documento también accesible en \texttt{/docs/json} y \texttt{/docs/yaml}; las rutas de \texttt{/v1/api} declaran esquemas y las rutas de \texttt{/v1/sim} los declaran de forma progresiva.

La Documentación Centralizada corresponde a un contenedor estático publicado mediante GitHub Pages y compuesto principalmente por archivos Markdown. Su propósito es reunir en un solo lugar la información necesaria para realizar pruebas con cada pasarela, incluyendo tarjetas de prueba, credenciales de sandbox, referencias de pago y tokens de acceso. La documentación está organizada por proveedor y tipo de dato, permitiendo que tanto el Desarrollador como el Tester puedan consultar la información desde un navegador mediante HTTPS y encontrar rápidamente los datos necesarios para el desarrollo y las pruebas.

\subsubsection{Relaciones entre contenedores}
En producción, el SDK se comunica directamente con las pasarelas de pago mediante HTTPS. Para las pruebas, las solicitudes se redirigen hacia la API de Simulación, también mediante HTTPS y utilizando JSON. La API de Simulación incluye la especificación OpenAPI 3.0 en el endpoint /docs, desde donde se pueden consultar los endpoints y estructuras disponibles. En el flujo productivo, las pasarelas pueden enviar eventos de forma asíncrona hacia la aplicación mediante webhooks sobre HTTPS.

\subsection{Diagrama de Contenedores — API de Simulación (Nivel 2)}
\begin{figure}[htbp]
\centering
\includegraphics[width=0.65\textwidth]{"../../architecture/Simulador API/Container Diagram - C4.png"}
\caption{Diagrama de Contenedores de la API de Simulación}
\label{fig:container_sim}
\end{figure}

Este diagrama complementa el anterior con un enfoque específico en el flujo de simulación, mostrando con mayor detalle cómo interactúan el SDK, la API de Simulación, la Documentación Interactiva y la Aplicación del Comercio durante las pruebas de desarrollo.

\subsubsection{Contenedores en el contexto de simulación}
El SDK Kit Pagos Colombia aparece en este diagrama como el cliente de la API de Simulación. Cuando se configura en modo simulación, sus adaptadores redirigen las solicitudes hacia el simulador en lugar de comunicarse con las APIs reales de los proveedores. Este cambio se realiza mediante la configuración del entorno, por lo que la aplicación que utiliza el SDK no necesita modificar su código para pasar de un entorno de simulación a uno de producción.

La API de Simulación Desplegada es el componente principal de este flujo. Recibe las solicitudes enviadas por los adaptadores del SDK y determina la respuesta que debe generar según el escenario configurado mediante los headers de control. A partir de esta información, la API genera una respuesta siguiendo el formato correspondiente a la pasarela que se está simulando. También permite generar webhooks sintéticos firmados hacia la Aplicación del Comercio cuando el Tester solicita explícitamente su envío. La API expone además la especificación OpenAPI 3.0 mediante el endpoint /docs. Desde allí, tanto el Desarrollador como el Tester pueden consultar los endpoints disponibles, los headers aceptados, los escenarios de simulación y ejemplos de respuesta. Esta documentación puede consultarse directamente desde el navegador o utilizarse como referencia para realizar pruebas desde herramientas como Postman. El servicio se encuentra desplegado en Render y cuenta con una URL pública mediante HTTPS, por lo que puede ser utilizado tanto por el SDK como directamente desde herramientas externas sin necesidad de realizar una configuración local del simulador.

La Aplicación del Comercio (Huésped) corresponde al sistema que integra el SDK para gestionar los pagos. Dentro del flujo de simulación, esta aplicación recibe los webhooks generados por la API de Simulación. Para esto, la aplicación debe disponer de un endpoint de notificaciones al cual el simulador pueda enviar los eventos mediante HTTP POST. Los eventos incluyen la firma correspondiente para representar el comportamiento que tendría una pasarela real al enviar una notificación en producción.

\subsubsection{Relaciones entre contenedores en el flujo de simulación}
El Desarrollador o Tester instala el SDK mediante npm y lo configura para trabajar en modo simulación. La Aplicación del Comercio utiliza las operaciones de pago expuestas por el SDK y, dependiendo de la configuración, los adaptadores internos envían las solicitudes hacia la API de Simulación mediante HTTPS y JSON. La API procesa cada solicitud de acuerdo con el escenario configurado y genera la respuesta correspondiente. Esta respuesta es enviada nuevamente al SDK, donde se normaliza para que la Aplicación del Comercio pueda trabajar con una entidad Transaction común, independientemente de la pasarela que se esté simulando.

Cuando se requiere probar el flujo de notificaciones, el Tester puede solicitar manualmente el envío de un webhook mediante el endpoint /v1/sim/webhooks/trigger. La API de Simulación genera entonces el evento correspondiente y lo envía mediante HTTP POST al endpoint configurado en la Aplicación del Comercio. Por otra parte, el Desarrollador o Tester puede consultar la documentación del simulador desde /docs. Esto permite revisar y probar los endpoints disponibles sin tener que consultar directamente el código fuente de la API.

